При 1225~К ширина Ti-богатой ветви остаётся 1,40~п.п. для 768 и 1400 атомов, а барьер увеличивается с 1,98 до 3,05~\(k_{\rm B}T\). Разрешение ветвей при 1240~К зависит от выбора разделителя. При \(q_c=0{,}47\) на 1400 атомах оценки ширины равны \(+0{,}30\) и \(-0{,}25\)~п.п., однако 95\%-й интервал максимума логарифма весов включает ноль.

На 2304 атомах перенос с 1245 на 1243~К прогнозирует \(+0{,}23\) и \(-0{,}21\)~п.п. на двух ветвях. Оба корня найдены в 53 из 64 блочных повторов. Их поля выходят за границы прямой сетки на 4,5 и 2,3~мэВ соответственно. При переносе на 1240~К эффективное число кадров порядка единицы.
